\begin{abstract}
We characterize the minimax separation rate for testing whether a conditional mean treatment effect is constant under a known uniform covariate distribution on the unit interval, fixed overlap, independent Hölder restrictions on the propensity, baseline mean, and effect, fixed mean envelopes, and an armwise conditional outcome-moment bound of order between one and two. Separation is measured by the population \(L_2\) distance of the effect from its average. Matching upper and lower bounds identify two rate mechanisms: effect approximation under outcome tails and the interaction of rough propensity and baseline functions. An explicit test combines outcome clipping, histogram projection corrections, independent evaluation blocks, and profiling over the unknown common effect. It controls size over the entire composite null at every stipulated sample size and attains uniform power at the stated separation whenever that separation lies below the model maximum, with public thresholds ensuring this condition. At matched smoothness, every moment order below two yields a polynomially larger critical radius than the bounded-outcome benchmark, whose minimax risk equals that of a signed-binary submodel. Supplying the entire true continuous propensity yields a matching frontier over a broader continuous-nuisance class retaining uniform design, overlap, effect smoothness, mean caps, and the conditional moment envelope. Comparing the two experiments on the original class identifies the exact region where the unknown-propensity experiment attains the supplied-propensity separation order.
\end{abstract}

\section{Introduction}

This paper establishes matching minimax separation bounds for conditional mean treatment-effect homogeneity under finite conditional outcome moments and independently specified nuisance smoothness. The statistical target is constancy of the difference between the treated and untreated conditional means, with an unknown common effect. Under a potential-outcome representation satisfying consistency and conditional independence of treatment assignment, this contrast is the conditional average treatment effect. We measure heterogeneity by its population \(L_2\) distance from its average, which equals the standard deviation of the conditional effect under the stipulated covariate distribution. The resulting frontier quantifies how outcome tails and assignment information jointly determine the smallest uniformly detectable departure from homogeneity.

The model has a known uniform covariate distribution on \([0,1]\), a propensity between \(1/4\) and \(3/4\), and separate Hölder restrictions on the propensity, baseline conditional mean, and mean effect. Write \(p\), the conditional moment order, for an exponent satisfying \(1<p\le2\); write \(\alpha\) and \(\beta\), the propensity and baseline smoothness exponents, for numbers in \((0,1]\); and write \(\gamma\), the effect smoothness exponent, for a number in \([1/4,1]\). These exponents are public. The Hölder radii are fixed at \(20\), the baseline mean and effect have supremum norms at most \(1/2\), and each treatment arm has conditional raw \(p\)-moment at most \(10\), uniformly in the covariate up to null sets. These restrictions define the class in \cref{obj:def:model}. Testing errors are assessed uniformly over that class: size covers every admissible constant effect and nuisance function, while power covers every law separated by the stipulated population distance, as specified in \cref{obj:def:testing-risk}.

Our main result gives the critical separation radius up to positive constants depending on the public tuple. Define \(q=(p-1)/p\), the moment-derived exponent, and \(S=\alpha+\beta\), the sum of nuisance smoothness exponents. The two separation exponents are
\[
E_0(p)=\frac{2\gamma q}{2\gamma+q},
\qquad
E_4(p)=
\frac{2S}{1+2\alpha+\beta/q+S/(2\gamma)}.
\]
For sample size \(n\ge2\), the capped critical radius has order
\[
n^{-\min\{E_0(p),E_4(p)\}},
\]
as established in \cref{obj:thm:heavy-tail-comparison}. The first exponent balances effect approximation with outcome-moment control. The second incorporates the interaction of propensity and baseline regularity. Their exact equality boundary is determined by
\[
F(p)=
\frac{2\alpha}{p-1}
+\frac{\beta}{q}
+\frac{S}{2\gamma},
\]
the phase-comparison quantity: effect smoothness determines the rate when \(F(p)\ge1\), and rough confounding determines it when \(F(p)<1\). The bounds are pure powers of sample size, including at equality, with constants locally uniform on compact subsets of the admissible exponent domain.

The attaining procedure operates directly on the original records. It clips outcomes, forms histogram scores with ordered-pair projection corrections, and takes the inner product of scores from two independent evaluation blocks. Profiling this quadratic statistic over the admissible common-effect interval accommodates the composite null. The public construction in \cref{obj:def:explicit-score-ledger} specifies all ranks, clipping levels, bias and covariance bounds, and rejection cutoffs. Its rejection probability is at most \(0.0075\) throughout the null for every \(n\ge2\). Its type-II error is at most \(0.0075\) at the stated upper separation when that separation lies below the model's maximal heterogeneity distance; fully evaluable sample-size thresholds ensure a nonempty power domain. These guarantees appear in \cref{obj:thm:whole-null-calibration-resolved,obj:thm:original-record-attainable-rate}. Matching converses apply to all admissible randomized tests observing the complete records.

The matched outcome comparison retains the same design, overlap, smoothness, and mean envelopes while imposing bounded outcomes. Its separation exponent is the original exponent evaluated at \(p=2\). Every admissible \(p<2\) strictly reduces that exponent, so the finite-moment critical radius exceeds its bounded counterpart by a polynomial factor. At \(p=2\), the ratio remains bounded across sample sizes. Moreover, mean-preserving sign randomization gives exact equality of bounded and signed-binary minimax risks and capped radii, as shown in \cref{obj:prop:bounded-submodel-comparison}. The full-record witnesses in \cref{obj:thm:tail-essential-full-record-converse} connect the slower finite-moment frontier to growing outcome marks: for each fixed tuple with \(p<2\), close finite-moment null and alternative mixtures coexist with eventual total-variation separation of every normalized finite prior pair supported on the bounded null and bounded alternatives at the same lower separation.

Supplying the entire true continuous propensity gives a complementary frontier. Over the continuous-carrier class in \cref{obj:def:oracle-broad-model}, the propensity and conditional arm means may have unrestricted moduli of continuity, while uniform design, overlap, effect Hölder smoothness, mean caps, and the conditional moment envelope are retained. A clipped inverse-propensity test attains the matching order \(n^{-E_0(p)}\) on this broader class, with whole-null calibration and a public nonempty-power threshold, by \cref{obj:thm:oracle-frontier-broad-class}. Comparing supplied and unknown propensity on the original shared class then yields an exact preservation boundary: their critical-radius ratio remains bounded over \(n\ge2\) precisely when \(F(p)\ge1\). For \(F(p)<1\), supplying the propensity improves separation by a factor of order \(n^{E_0(p)-E_4(p)}\). This comparison, including its equality case, follows from the matched frontiers in \cref{obj:thm:oracle-frontier-resolved,obj:thm:exact-propensity-preservation-boundary}.

These results connect econometric homogeneity testing with quadratic nonparametric testing and higher-order causal methods. \citet[Section III]{CrumpHotzImbensMitnik2008Homogeneity} develop sieve Wald tests of conditional mean-effect homogeneity. Recent orthogonal-score procedures provide asymptotic calibration and power under their respective outcome and nuisance-learning conditions \citep[Assumptions 1--5, Theorems 3.1--3.3 and 4.1]{LuSong2026OrthogonalICM}; \citep[Assumptions 3.3--3.4, Proposition 3.1]{LapentaStrittmatterVergara2026Omnibus}. Directional testing also explicitly pursues uniform asymptotic calibration, with power tied to alignment of a learned direction and the population signal \citep[Theorem 1, Corollaries 1--2]{DhawanGuoShah2026Directional}. Our contribution concerns finite-sample whole-null calibration and uniform population separation over the primitive class specified above. The approximation--variability balance has antecedents in quadratic goodness-of-fit theory \citep[Sections 2--4, Theorems 1--2]{IngsterSapatinas2009GoodnessOfFit}, while projection corrections connect to higher-order conditional-effect estimation \citep[Sections 3--4]{KennedyBalakrishnanRobinsWasserman2024MinimaxCATE}.



The paper proceeds through related work, setup and assumptions, the matched separation rates and outcome-tail comparison, the calibrated original-record construction, and the supplied-propensity frontiers and assignment-information comparison. A discussion considers interpretation and extensions. The appendices develop causal identification and benchmark equivalence, calibration and attainment, full-record lower bounds, supplied-propensity arguments and comparison algebra, followed by proofs of the main results.
